% Part: first-order-logic
% Chapter: completeness
% Section: introduction

\documentclass[../../../include/open-logic-section]{subfiles}

\begin{document}

\iftag{FOL}
      {\olfileid{fol}{com}{int}}
      {\olfileid{pl}{com}{int}}

\olsection{પરિચય}

પૂર્ણતા પ્રમેય તર્કશાસ્ત્ર અંગેના સૌથી પાયાના પરિણામોમાંનું એક છે.
તેના બે નિરૂપણ છે, અને એ બંને સમતુલ્ય છે તે આપણે સાબિત કરીશું. તેના
પહેલા નિરૂપણમાં તે અર્થાનુસારી ફલિતતા અને આપણા !!{derivation} તંત્ર
વચ્ચેના સંબંધ અંગે પાયાની વાત કહે છે: જો અમુક !!{sentence}s~$\Gamma$માંથી
!!a{sentence}~$!A$ ફલિત થતું હોય, તો $\Gamma \Proves !A$ સ્થાપિત કરતી
!!a{derivation} પણ હોય છે. આમ, ખરેખર ફલિત ન થતી બાબતો સાબિત કર્યા વિના
!!{derivation} તંત્ર શક્ય હોય એટલું સમર્થ છે.

તેનું બીજું નિરૂપણ સંરચના-અસ્તિત્વના પરિણામ તરીકે આપી શકાય છે:
!!{sentence}sનો દરેક સુસંગત ગણ સંતોષ્ય છે. સુસંગતતા સાબિતી-સૈદ્ધાંતિક
ખ્યાલ છે: તે કહે છે કે આપણું !!{derivation} તંત્ર અમુક પ્રકારની
!!{derivation}s બનાવી શકતું નથી. પરંતુ~$\Gamma$માંથી અમુક પ્રકારની
!!{derivation}s નથી એટલા પરથી જ $\iftag{FOL}{\Sat{M}}{\pSat{v}}{\Gamma}$
હોય એવી \iftag{FOL}{!!a{structure}~$\Struct{M}$}{!!{valuation}{~$\pAssign{v}$}}
હોવાની ખાતરી મળે છે એમ કોણ કહે? પૂર્ણતા પ્રમેય પહેલી વાર સાબિત થયો
તે પહેલાં—વાસ્તવમાં, અત્યારે આપણી પાસે છે તેવાં !!{derivation} તંત્રો
મળ્યાં તે પહેલાં—મહાન જર્મન ગણિતશાસ્ત્રી ડેવિડ હિલ્બર્ટ એવું માનતા કે
ગાણિતિક સિદ્ધાંતોની સુસંગતતા તેઓ જેના વિશે છે તે પદાર્થોના અસ્તિત્વની
ખાતરી આપે છે. ગૉટલૉબ ફ્રેગેને લખેલા પત્રમાં તેમણે આ વાત આ રીતે મૂકી:
\begin{quote}
  સ્વેચ્છાએ આપેલી સ્વયંસિદ્ધિઓ તેમનાં બધાં પરિણામો સહિત પરસ્પર
  વિરોધાભાસી ન હોય, તો તે સાચી છે અને સ્વયંસિદ્ધિઓથી વ્યાખ્યાયિત
  વસ્તુઓ અસ્તિત્વ ધરાવે છે. મારા માટે સત્ય અને અસ્તિત્વની આ કસોટી છે.
\end{quote}
ફ્રેગે તેનો ઉગ્ર વિરોધ કર્યો. પૂર્ણતા પ્રમેયનું બીજું નિરૂપણ બતાવે છે
કે ઓછામાં ઓછા એટલા અર્થમાં હિલ્બર્ટ સાચા હતા કે સ્વયંસિદ્ધિઓ સુસંગત
હોય, તો તેમને બધાને સાચી બનાવતી \emph{કોઈક}
\iftag{FOL}{!!{structure}}{!!{valuation}} અસ્તિત્વ ધરાવે છે.

પૂર્ણતા પ્રમેય—અથવા વધુ ચોક્કસ રીતે, તેની સાબિતી—મહત્વની હોવાનું આટલું
જ કારણ નથી. તેનાં અનેક મહત્વનાં પરિણામો છે, જેમાંનાં કેટલાંકની આપણે
અલગથી ચર્ચા કરીશું. દાખલા તરીકે, $\Gamma \Proves !A$ બતાવતી કોઈ પણ
!!{derivation} સાન્ત છે અને તેથી~$\Gamma$નાં સાન્ત સંખ્યામાં
!!{sentence}s જ વાપરી શકે છે. આથી પૂર્ણતા પ્રમેય મુજબ જો~$!A$ એ
~$\Gamma$નું પરિણામ હોય, તો તે~$\Gamma$ના કોઈ સાન્ત ઉપગણનું પરિણામ
પહેલેથી જ હોય છે. તેને \emph{સઘનતા} કહે છે. સમતુલ્ય રીતે, જો
$\Gamma$નો દરેક સાન્ત ઉપગણ સુસંગત હોય, તો~$\Gamma$ પોતે પણ સુસંગત
હોવો જ જોઈએ.

સઘનતા પ્રમેય !!{derivation}s મારફતના ફેરા વડે પૂર્ણતા પ્રમેયમાંથી
ફલિત થાય છે, છતાં પૂર્ણતા પ્રમેયની \emph{સાબિતી} વાપરીને તેને સીધો
સ્થાપિત કરવો પણ શક્ય છે. સાબિતી કોઈ ચોક્કસ ગુણધર્મ—સુસંગતતા—ધરાવતા
!!{sentence}sનો ગણ લે છે અને તેમાંથી ચોક્કસ ગુણધર્મો ધરાવતી
!!a{structure} રચે છે (આ કિસ્સામાં, તે ગણને સંતોષે છે). સુસંગત ગણોને
બદલે !!{sentence}sના ``સાન્ત રીતે સંતોષ્ય'' ગણોથી શરૂ કરીને લગભગ આ જ
રચના વડે સઘનતા સીધી સ્થાપિત કરી શકાય છે.
\sourcecorrection{OLCO-001}{મૂળની \texttt{introduction.tex},
પંક્તિ~61માં ``the proof of'' પહેલાં અનાવશ્યક બીજો ``the'' આવે છે.
અહીં એ પુનરુક્તિ દૂર કરી છે.}
\iftag{FOL}{આ રચના બીજાં પરિણામો પણ આપે છે; દાખલા તરીકે,
  !!{sentence}sના દરેક સંતોષ્ય ગણને સાન્ત અથવા !!a{denumerable}
  સંરચના મળે છે. (આ પરિણામને લેવેનહાઇમ--સ્કોલેમ પ્રમેય કહે છે.)
  સામાન્ય રીતે, !!{sentence}sના ગણોમાંથી !!{structure}s રચવાની રીત
  તર્કશાસ્ત્રમાં વારંવાર અને ક્યારેક તો તત્ત્વજ્ઞાનમાં પણ વપરાય છે.}{}
\sourcecorrection{OLCO-002}{મૂળની \texttt{introduction.tex},
પંક્તિ~70માં વિશેષણરૂપ \texttt{denumerable} પછી વધારાનો
\texttt{s} છે. અહીં તેને એકવચન ``સંરચના'' સાથે મેળ ખાતા વિશેષણરૂપે
લખ્યું છે.}

\end{document}
